\nextpage
\section{Как проводились измерения}
Для каждого входа независимая программа на Python заранее выписала ожидаемые тройки «вид, начало, конец». До прогрева стенд сравнил с ними каждую тройку всех четырёх сканеров, включая EOF. После EOF проверен повторный вызов. Отдельные тесты охватили все 31 вида, пустой текст, ошибочный байт, комментарий без перевода строки и лексему у конца буфера.

\subhead{Что означает наносекунда на байт}
Пусть $N$~--- число байтов в файле, включая пробелы, комментарии и переводы строк. $K$~--- число выданных обычных лексем, без EOF. Перед запуском файл уже находился в памяти и экземпляры сканера были открыты. Счётчик времени включали перед первым \code{next}. В интервал вошли вызовы до EOF включительно, подсчёт лексем и одинаковое для четырёх вариантов обновление 64-битной контрольной суммы по виду и обеим границам. Закрытие экземпляров, сверку с эталоном и запись результата выполняли после остановки счётчика.

Пусть $E$~--- длительность измеренного интервала в наносекундах, $P$~--- число полных проходов в нём. При $N<65\,536$ выполняли $P=16$ проходов, иначе $P=1$. Время одного прохода $T=E/P$. В таблицы вошёл показатель
\[
             t=\frac{T}{N}=\frac{E}{P\,N}\quad\text{нс/байт}.
\]
Например, для корпуса $P\,N=16\cdot17\,430=278\,880$ байт. Так учитываются и байты, которые сканер прочитал и пропустил. Показатель характеризует полную обработку входа вместе с выдачей результата.

\begin{center}\captionof{table}{Условия общей серии А и Б.}\label{tab:env}
\small\begin{tabularx}{\linewidth}{@{}lX@{}}\toprule
Процессор и система & Intel Core i5-13400F, Windows 11 25H2, сборка 26200\\
Компиляция & Компиляторы \code{gcc} и \code{g++} 13.2.0 из MinGW-w64. Параметры оптимизации \code{-O3 -DNDEBUG}. Стенд и сгенерированные сканеры написаны на C11, собственный сканер~--- на C++17.\\
Генераторы и часы & win\_flex 2.6.4 (\code{-8 -Cf}), re2c 4.6. Время измерено системным счётчиком \code{QueryPerformanceCounter}.\\
Порядок & Один прогрев на вариант и вход, затем 16 раундов. Первый порядок: flex, re2c, table, branch. В каждом следующем раунде первый вариант перемещался в конец.\\ \bottomrule
\end{tabularx}\end{center}
Для каждого входа и варианта указаны медиана и квартили $Q_1$, $Q_3$ по 16 раундам. Для доли $p$ берётся позиция $15p$ в отсортированном ряду с нумерацией от нуля и линейной интерполяцией. Интервал между квартилями описывает среднюю половину замеров.
